\cotefile{42}
\scanpage{2}
\uline{Réseau\unsure{’|} d'anneaux locaux}

est formé d'un ensemble ordonné $X$, et \deleted{pour \unsure{tout|tt} $x\in X$, d} d'un système in\uline{du}ctif d'anneaux locaux $(O_x)_{x\in X}$, [donc si $x\leq y$, on a $O_x\to O_y$) tel que pour $x\leq y$, $O_x\to O_y$ induise un isomorphisme $(\underline{O}_x)_{p_{xy}}\to\underline{O}_y$, (où $p_{xy}$ est l'image \unsure{inverse|réciproque} de $\unsure{r|\mathfrak{r}|\mathfrak{m}}(\underline{O}_y)$) et que \deleted{$p_{xy}\unsure{\supset|\supseteq}p_{xy'}\Rightarrow y\unsure{\leq|\geq}y'$} \inserted{l'on obtienne ainsi une correspondance biunivoque entre les $y\geq x$ et les idéaux premiers de $O_x$}.
% slanted
\margin{Notation $X,Y$ pour des réseaux.}

Un hom. d'un réseau $X$ dans un autre $Y$ : hom des systèmes inductifs d'anneaux \uline{locaux}, i.e. \unsure{formé} d'une application \deleted{\unsure{continue}} croissante $X\xrightarrow{f}Y$, et pour \unsure{tt|tout} $x\in X$ d'un hom \uline{local} $\underline{O}_{f(x)}\to\underline{O}_x$, compatible avec les applications de \deleted{\unsure{spécialisati|spécialisation}} générisation.

Soit $\unsure{\mathcal R|\mathscr R|R}$ la catégorie des réseaux, $\unsure{S|\mathcal S}$ la catégorie des préschémas. On a un foncteur canonique $S\xrightarrow{\unsure{\rho|P|p}}\mathcal R$.

Réseau sur un autre, sur un préschéma.

Soit $S$ un préschéma \unsure{loc|loc.} noethérien,
% archive sheet number: 1
